\chapter{Asymmetric scale doubling and rational certificates}
\label{chap:chain}

Theorem~\ref{thm:gap-of-seeds} requires both initial purity defects to be at most \(1/100\). On
moderate-size rings, certified seed defects are typically larger and unequal, such as
\((\es, \et) \approx (0.064, 0.041)\), and enter the \(\le 1/100\) regime only after several
scale-doubling steps. To cover such seeds, we iterate the asymmetric scale-doubling recurrence for
finitely many steps until both defects fall below \(1/100\), after which
Theorem~\ref{thm:gap-of-seeds} applies at all larger scales. Because the recurrence maps are
rational functions, every finite-step inequality can be verified in exact rational arithmetic.

\section{The asymmetric scale-doubling recurrence}

Rather than bounding both defects by a single common value \(\varepsilon\), the asymmetric
recurrence tracks separate upper bounds \(s\) and \(t\) on the spatial defect
\(1 - \Ps(n, \beta)\) and the thermal defect \(1 - \Pt(2n, \beta)\) through the coupled defect
map~\eqref{eq:coupling-defect} (Definition~\ref{def:coupling-defect}) and the purity-squaring
map~\eqref{eq:squaring-map} (Definition~\ref{def:squaring-defect}). Because the thermal half of
each step (Lemma~\ref{lem:doubling-step}) uses the already-squared spatial defect
\(s' \ge f(z(s, t))\) rather than the initial spatial defect \(s\), this decoupled update has a
substantially larger basin of attraction than the symmetric recurrence of Lemma~\ref{lem:doubling}
(see~\cite{ScopeNote}, Section~E).

\begin{definition}[Coupled defect map]
  \label{def:coupling-defect}
  \lean{CertifiedGaps.couplingDefect}
  \leanok
  For defect bounds \(a, b \in \R\), define
  \begin{equation}
    \label{eq:coupling-defect}
    z(a, b) = 1 - (1 - b)(1 - a)^2.
  \end{equation}
  Whenever two purities satisfy \(P_1 \ge 1 - a\) and \(P_2 \ge 1 - b\), the product \(P_1^2 P_2\)
  appearing in the spacetime coupling identities~\eqref{eq:coupling} of Lemma~\ref{lem:coupling} is
  at least \(1 - z(a, b)\).
\end{definition}

\begin{lemma}[Monotonicity of the spectral gap bound]
  \label{lem:gapgt-mono}
  \lean{CertifiedGaps.GapGT.mono}
  \leanok
  \uses{def:gapgt}
  If \(e\) has a spectral gap greater than \(c\) (Definition~\ref{def:gapgt}) and \(c' \le c\), then
  \(e\) has a spectral gap greater than \(c'\).
\end{lemma}
\begin{proof}
  \leanok
  The same minimizing index \(m_0\) in Definition~\ref{def:gapgt} witnesses the bound for \(c'\).
\end{proof}

\begin{lemma}[Defect form of purity squaring]
  \label{lem:defect-squaring}
  \lean{CertifiedGaps.one_sub_thermalPurity_two_mul_le,
    CertifiedGaps.one_sub_spatialPurity_two_mul_le, CertifiedGaps.squaringDefect_le_of_le}
  \leanok
  \uses{lem:power-sums, lem:squaring-homogeneous, lem:f-mono, def:moment-bounds}
  For any \(L\) and \(\beta\), \(1 - \Pt(L, 2\beta) \le f\bigl(1 - \Pt(L, \beta)\bigr)\), and under
  the spatial moment inequalities of Definition~\ref{def:moment-bounds} above \(n_0\),
  \(1 - \Ps(2n, \beta) \le f\bigl(1 - \Ps(n, \beta)\bigr)\) for even \(n \ge n_0\) and
  \(\beta > 0\). If \(P \le 1\) and \(1 - z \le P\) with \(z < 1\), then \(f(1 - P) \le f(z)\).
\end{lemma}
\begin{proof}
  \leanok
  Apply Lemma~\ref{lem:squaring-homogeneous} to the Boltzmann weights (via~\eqref{eq:power-sums} of
  Lemma~\ref{lem:power-sums}) for the thermal purity and to condition (2) of
  Definition~\ref{def:moment-bounds} for the spatial purity; the final inequality follows from the
  monotonicity of \(f\) in Lemma~\ref{lem:f-mono}.
\end{proof}

\begin{lemma}[Asymmetric scale-doubling step]
  \label{lem:doubling-step}
  \lean{CertifiedGaps.doubling_step}
  \leanok
  \uses{def:coupling-defect, lem:coupling, lem:defect-squaring, lem:purity-basic,
    lem:purity-squaring}
  Assume the spatial moment inequalities of Definition~\ref{def:moment-bounds} above \(n_0\), let
  \(n \ge n_0\) be even, and let \(\beta > 0\). Let \(s, t, s' \le 1\) satisfy
  \begin{equation}
    \label{eq:asym-step}
    z(s, t) < 1, \qquad f\bigl(z(s, t)\bigr) \le s', \qquad
    z(t, s') < 1, \qquad f\bigl(z(t, s')\bigr) \le t'.
  \end{equation}
  If \(\Ps(n, \beta) \ge 1 - s\) and \(\Pt(2n, \beta) \ge 1 - t\), then
  \(\Ps(2n, 2\beta) \ge 1 - s'\) and \(\Pt(4n, 2\beta) \ge 1 - t'\).
\end{lemma}
\begin{proof}
  \leanok
  Combining the first spacetime coupling identity in~\eqref{eq:coupling} (Lemma~\ref{lem:coupling})
  with \(\Pt(n,\beta) \le 1\) (Lemma~\ref{lem:purity-basic}) and~\eqref{eq:coupling-defect} gives
  \(\Ps(n, 2\beta) \ge 1 - z(s, t)\), and spatial purity squaring (Lemma~\ref{lem:defect-squaring})
  together with~\eqref{eq:asym-step} bounds \(\Ps(2n, 2\beta) \ge 1 - s'\). Substituting that
  updated spatial bound into the second coupling identity in~\eqref{eq:coupling} along with
  \(\Ps(2n, \beta) \le 1\) (Lemma~\ref{lem:purity-squaring}) gives
  \(\Pt(4n, \beta) \ge 1 - z(t, s')\), and thermal purity squaring (Lemma~\ref{lem:defect-squaring})
  bounds \(\Pt(4n, 2\beta) \ge 1 - t'\).
\end{proof}

\begin{lemma}[Dyadic length interpolation]
  \label{lem:window}
  \lean{CertifiedGaps.thermalPurity_ge_of_window}
  \leanok
  \uses{def:coupling-defect, def:moment-bounds, lem:purity-basic}
  Under the spatial moment inequalities of Definition~\ref{def:moment-bounds} above \(n_0\), let
  \(n \ge n_0\) be even and positive, \(\beta > 0\), \(s \le 1\), \(\Ps(n, \beta) \ge 1 - s\), and
  \(\Pt(2n, \beta) \ge 1 - t\). Then every even \(L\) with \(n \le L \le 2n\) satisfies
  \(\Pt(L, \beta) \ge 1 - z(s, t)\).
\end{lemma}
\begin{proof}
  \leanok
  Repeating the \(\ell^p\)-interpolation estimate~\eqref{eq:dyadic-interp} from the proof of
  Theorem~\ref{thm:gap-of-seeds} with~\eqref{eq:coupling-defect} gives
  \(\Pt(L, \beta) \ge x^{L/(2n)} \ge x = \Pt(2n, \beta)\, \Ps(n, \beta)^2 \ge 1 - z(s, t)\).
\end{proof}

\begin{lemma}[Spectral gap from a thermal defect bound]
  \label{lem:thermal-gap}
  \lean{CertifiedGaps.gapGT_of_thermalPurity}
  \leanok
  \uses{lem:purity-to-gap}
  If \(\beta > 0\), \(z < 1\), \(\Pt(L, \beta) \ge 1 - z\), and \(e^{\beta c} z < 1 - z\), then
  \(E_L\) has a spectral gap greater than \(c\).
\end{lemma}
\begin{proof}
  \leanok
  At any excited level \(m \ne m_0\),~\eqref{eq:purity-to-gap} in Lemma~\ref{lem:purity-to-gap}
  gives \(e^{-\beta (E_L(m) - E_L(m_0))} \le (1 - \Pt)/\Pt \le z/(1-z) < e^{-\beta c}\); taking
  logarithms and dividing by \(\beta\) yields \(E_L(m) - E_L(m_0) > c\).
\end{proof}

\section{Multi-step asymmetric doubling}

A finite sequence of defect pairs \((s_j, t_j)_{j \le J}\) satisfying~\eqref{eq:asym-step} connects
the initial seed at \(j = 0\) to scale \(J\) (Lemma~\ref{lem:chain-purities}), where both defects
satisfy \(s_J, t_J \le \varepsilon \le 1/100\). On each initial dyadic interval
\([n_0 2^j, 2n_0 2^j]\) with \(j < J\), Lemmas~\ref{lem:window} and~\ref{lem:thermal-gap} bound the
spectral gap directly; for all lengths \(L \ge n_0 2^J\), Theorem~\ref{thm:gap-of-seeds} applies
starting from the rescaled seed \((n_0 2^J, \beta_0 2^J)\) (Theorem~\ref{thm:gap-of-chain}).

\begin{lemma}[Iterated asymmetric doubling]
  \label{lem:chain-purities}
  \lean{CertifiedGaps.purities_ge_chain}
  \leanok
  \uses{lem:doubling-step}
  Assume the spatial moment inequalities of Definition~\ref{def:moment-bounds} above an even \(n_0\)
  and \(\beta_0 > 0\). Let \((s_j, t_j)_{j \le J}\) satisfy \(s_j, t_j \le 1\) and, for \(j < J\),
  the step conditions~\eqref{eq:asym-step} of Lemma~\ref{lem:doubling-step} with
  \((s, t, s', t') = (s_j, t_j, s_{j+1}, t_{j+1})\). If \(\Ps(n_0, \beta_0) \ge 1 - s_0\) and
  \(\Pt(2n_0, \beta_0) \ge 1 - t_0\), then for all \(j \le J\),
  \(\Ps(n_0 2^j, \beta_0 2^j) \ge 1 - s_j\) and \(\Pt(2 n_0 2^j, \beta_0 2^j) \ge 1 - t_j\).
\end{lemma}
\begin{proof}
  \leanok
  Induct on \(j\) using Lemma~\ref{lem:doubling-step}.
\end{proof}

\begin{theorem}[Asymmetric spacetime-bootstrap criterion]
  \label{thm:gap-of-chain}
  \lean{CertifiedGaps.gapGT_of_chain}
  \leanok
  \uses{lem:chain-purities, lem:dyadic-window, lem:window, lem:thermal-gap, thm:gap-of-seeds,
    lem:gapgt-mono, lem:moment-bounds-limit}
  Assume the spatial moment inequalities of Definition~\ref{def:moment-bounds} above an even
  \(n_0 > 0\), \(\beta_0 > 0\), seed bounds \(\Ps(n_0, \beta_0) \ge 1 - s_0\),
  \(\Pt(2n_0, \beta_0) \ge 1 - t_0\), and a defect sequence \((s_j, t_j)_{j \le J}\) as in
  Lemma~\ref{lem:chain-purities}. Let \(r > 0\) and \(0 < \varepsilon \le 1/100\) satisfy
  \begin{equation}
    \label{eq:chain-gap-cond}
    r^{2^j} z(s_j, t_j) < 1 - z(s_j, t_j) \quad (j < J), \qquad
    s_J, t_J \le \varepsilon, \qquad r^{2^J} \cdot 5\varepsilon \le 1.
  \end{equation}
  Then \(E_L\) has a spectral gap greater than \(\log r / \beta_0\) for every even \(L \ge n_0\).
\end{theorem}
\begin{proof}
  \leanok
  By Lemma~\ref{lem:dyadic-window}, choose \(j \in \N\) with \(n_0 2^j \le L \le 2 n_0 2^j\). When
  \(j < J\), Lemma~\ref{lem:window} and Lemma~\ref{lem:chain-purities} give
  \(\Pt(L, \beta_0 2^j) \ge 1 - z(s_j, t_j)\), and Lemma~\ref{lem:thermal-gap} together with the
  first condition in~\eqref{eq:chain-gap-cond} yields the gap bound. When \(j \ge J\),
  Theorem~\ref{thm:gap-of-seeds} applied at the rescaled seed \((n_0 2^J, \beta_0 2^J)\) with
  defect \(\varepsilon\) gives the lower bound~\eqref{eq:sym-gap-bound}, namely
  \(\log(1/(5\varepsilon)) / (\beta_0 2^J)\), which is at least \(\log r / \beta_0\) by the third
  condition in~\eqref{eq:chain-gap-cond}; Lemma~\ref{lem:gapgt-mono} then completes the proof.
\end{proof}

\section{Rational doubling certificates}

Because \(z(a, b)\) in~\eqref{eq:coupling-defect} and \(f(z)\) in~\eqref{eq:squaring-map} are
rational functions, packaging \((s_j, t_j)\), \(r\), and \(\varepsilon\) as rational numbers reduces
the hypotheses~\eqref{eq:asym-step} and~\eqref{eq:chain-gap-cond} of Theorem~\ref{thm:gap-of-chain}
to finitely many decidable rational inequalities (Definition~\ref{def:chain-certificate} and
Theorem~\ref{thm:gap-of-chain-certificate}).

\begin{definition}[Rational doubling certificate]
  \label{def:chain-certificate}
  \lean{CertifiedGaps.ChainCertificate, CertifiedGaps.ChainCertificate.Valid}
  \leanok
  \uses{def:coupling-defect, def:squaring-defect}
  A \emph{rational doubling certificate} consists of rational sequences \(s = (s_0, \dots, s_J)\)
  and \(t = (t_0, \dots, t_J)\), a decay base \(r \in \Q\), and a terminal defect bound
  \(\varepsilon \in \Q\). It is \emph{admissible} if \(r > 0\), \(0 < \varepsilon \le 1/100\),
  \(s_j, t_j \le 1\) for \(j \le J\), and both the recurrence inequalities~\eqref{eq:asym-step} for
  \(j < J\) and the terminal conditions~\eqref{eq:chain-gap-cond} hold in \(\Q\).
\end{definition}

\begin{lemma}[Compatibility with real arithmetic]
  \label{lem:chain-cast}
  \lean{CertifiedGaps.ChainCertificate.cast_zQ, CertifiedGaps.ChainCertificate.cast_fQ}
  \leanok
  \uses{def:chain-certificate}
  The rational coupled-defect map~\eqref{eq:coupling-defect} and purity-squaring
  map~\eqref{eq:squaring-map} agree with their real counterparts under the field inclusion
  \(\Q \hookrightarrow \R\).
\end{lemma}
\begin{proof}
  \leanok
  Both maps are given by the same rational expressions over the two fields, and the inclusion of
  the rationals into the reals commutes with the field operations, division included.
\end{proof}

\begin{theorem}[Spectral gap from a rational doubling certificate]
  \label{thm:gap-of-chain-certificate}
  \lean{CertifiedGaps.gapGT_of_chainCertificate}
  \leanok
  \uses{thm:gap-of-chain, def:chain-certificate, lem:chain-cast}
  Assume the spatial moment inequalities of Definition~\ref{def:moment-bounds} above an even
  \(n_0 > 0\), \(\beta_0 > 0\), an admissible rational doubling certificate
  (Definition~\ref{def:chain-certificate}) with initial defects \((s_0, t_0)\) and decay base \(r\),
  and seed bounds \(\Ps(n_0, \beta_0) \ge 1 - s_0\), \(\Pt(2n_0, \beta_0) \ge 1 - t_0\). Then
  \(E_L\) has a spectral gap greater than \(\log r / \beta_0\) for every even \(L \ge n_0\).
\end{theorem}
\begin{proof}
  \leanok
  Embed the rational inequalities of Definition~\ref{def:chain-certificate} into \(\R\) via
  Lemma~\ref{lem:chain-cast} and apply Theorem~\ref{thm:gap-of-chain}.
\end{proof}

\begin{remark}
  The rational entries of each doubling certificate are precomputed from the certified seed
  defects; this external computation need not be trusted because every inequality of
  Definition~\ref{def:chain-certificate} is verified directly in exact rational arithmetic.
\end{remark}
